\subsection{The four quality dimensions, and least privilege}

The four failures that open this paper come from two published taxonomies: completeness and
plausibility from Kahn et al.\ (2016), consistency and timeliness from Weiskopf and Weng's (2013)
concordance and currency. Consistency is raised to top level for cross-record checks. Conformance
is skipped, staying within one lineage since Weiskopf and Weng are among Kahn's co-authors.
Completeness carries the word both taxonomies leave implicit, missingness: Wells et al.\ (2013) find
that missing values in EHR-derived data are seldom missing at random. An absent field is itself
information, and the strategy chosen for handling it can change what the data support. None of this
says what a checker must \emph{see}: a plausibility check needs code, value and unit, hence
whole-record access by default.

Saltzer and Schroeder (1975) ask the opposite, minimal privilege for every program and user, a
concern with error as much as malice (Section~\ref{sec:threat}). Three mechanisms grant it, each in a
unit larger than the field.

\begin{itemize}[leftmargin=1.4em, itemsep=0.15em]
  \item \textbf{Saltzer and Schroeder's own principle.} The unit is the principal.
  \item \textbf{The HIPAA minimum necessary standard.} The unit is the class of worker.
  \item \textbf{SMART on FHIR scopes.} The unit is the resource type, as in
        \texttt{patient/Observation.read} (Mandel et al., 2016; HL7 International, 2024).
\end{itemize}

None
governs what one component hands another inside a pipeline, so a perimeter-compliant service can
still hand every field to the checks behind it, as the full-access \textbf{Baseline} of
Section~\ref{sec:designs} does.
